\chapter{Laparoscopic sterilisation}
\index{Laparoscopic sterilisation}

\section*{Definition and Overview}
Laparoscopic sterilisation is a permanent surgical method of female contraception that involves the occlusion or removal of the fallopian tubes via keyhole surgery (laparoscopy). By blocking or removing the pathways through which ova travel from the ovaries to the uterus, the procedure prevents sperm from fertilising an egg. It is a highly effective, irreversible procedure intended for women who are certain they do not wish to have children in the future. As a minimally invasive gynaecological intervention, it offers advantages such as rapid recovery and minimal scarring compared to open surgical approaches, though it requires general anaesthesia and carries specific surgical and post-operative risks.

\section*{Epidemiology}
Voluntary female sterilisation is one of the most widely used methods of contraception globally, particularly among women who have completed their families. In many countries and other developed countries, the prevalence of laparoscopic sterilisation has decreased slightly in recent decades due to the rising popularity and accessibility of long-acting reversible contraceptives (LARCs), such as hormonal and copper intrauterine devices (IUDs) and subdermal implants, as well as vasectomy. Nevertheless, it remains a common choice for women in their late 30s and 40s. The procedure has no significant age or racial predisposition beyond reproductive maturity and personal contraceptive choice, but rates of post-operative regret are statistically higher among individuals sterilised under the age of 30 or during times of personal crisis.

\section*{Aetiology and Pathophysiology}
Laparoscopic sterilisation is indicated for individuals seeking permanent contraception, or where pregnancy poses a severe, life-threatening medical risk. The underlying mechanism and surgical physiology include:
\begin{itemize}
    \item \textbf{Anatomical disruption:} The physical blockage or removal of the fallopian tubes prevents the fusion of sperm and mature oocytes within the ampulla.
    \item \textbf{Laparoscopic access:} Carbon dioxide gas is insufflated into the peritoneal cavity (pneumoperitoneum) to create a working space, permitting high-definition visualization of the pelvic organs using a laparoscope.
    \item \textbf{Mechanical occlusion:} Application of mechanical devices, such as titanium clips (e.g. Filshie clips) or silicone bands, to compress and shut the lumen of the mid-isthmic portion of the fallopian tubes.
    \item \textbf{Thermal occlusion:} Use of bipolar electrocautery to desiccate and destroy a segment of the tube, leading to subsequent fibrosis and closure.
    \item \textbf{Salpingectomy:} Complete excision of both fallopian tubes (bilateral salpingectomy), which is increasingly preferred over occlusion as it eliminates the risk of tubal pregnancy and significantly reduces the lifetime risk of high-grade serous ovarian cancer, which often originates in the fimbria of the tubes.
\end{itemize}

\section*{Clinical Features}
The clinical presentation of a patient undergoing laparoscopic sterilisation includes pre-operative planning and expected post-operative symptoms. Key features include:
\begin{itemize}
    \item \textbf{Elective presentation:} Patients present electively for a planned day-surgery procedure under general anaesthesia.
    \item \textbf{Incisions:} Typically, two small abdominal incisions are present: one umbilical incision (5--10 mm) for the laparoscope, and one suprapubic incision (5 mm) for the surgical instruments.
    \item \textbf{Post-operative abdominal discomfort:} Mild to moderate pelvic or abdominal cramping, similar to menstrual pain, in the first 24--48 hours.
    \item \textbf{Referred shoulder tip pain:} Diaphragmatic irritation caused by residual carbon dioxide gas trapped in the peritoneal cavity, which stimulates the phrenic nerve and projects pain to the shoulders.
    \item \textbf{Vaginal spotting:} Mild vaginal bleeding or spotting for a few days post-procedure, resulting from uterine manipulation during surgery.
    \item \textbf{Preserved endocrine function:} Ovarian function, hormone production (oestrogen and progesterone), and the menstrual cycle remain entirely unchanged, with menopause occurring at the biologically natural time.
\end{itemize}

\section*{Differential Diagnosis}
When evaluating post-operative concerns following laparoscopic sterilisation, clinicians must distinguish expected recovery symptoms from serious complications:
\begin{itemize}
    \item \textbf{Pelvic inflammatory disease (PID) or peritonitis:} Severe, constant abdominal pain accompanied by high fever and purulent vaginal discharge, distinguishing it from routine post-operative cramping.
    \item \textbf{Internal haemorrhage:} Persistent tachycardia, hypotension, pallor, and progressive abdominal distension or pain, which must be differentiated from benign incisional discomfort.
    \item \textbf{Visceral injury:} Severe, worsening abdominal pain, persistent vomiting, and abdominal rigidity arising from accidental thermal or mechanical damage to the bowel or bladder during trocar insertion.
    \item \textbf{Wound infection:} Localised erythema, warmth, swelling, and purulent drainage at the incision sites, distinct from normal wound healing.
    \item \textbf{Ectopic pregnancy:} In the rare event of late sterilisation failure, any presentation of pelvic pain and vaginal bleeding must be investigated for a tubal pregnancy rather than a normal menstrual cycle or pelvic pain of other origins.
\end{itemize}

\section*{When to Seek Medical Care}
Patients should monitor their recovery closely and contact their gynaecologist, general practitioner, or hospital emergency department if they experience progressive or unexpected symptoms, such as worsening pelvic pain, persistent vomiting, high fever, or spreading redness around the wounds.

\textbf{Contact your local emergency services for:}
\begin{itemize}
    \item Severe, sudden, or crushing chest pain.
    \item Sudden shortness of breath or difficulty breathing.
    \item Loss of consciousness, severe dizziness, or confusion.
    \item Signs of severe shock, such as cold, clammy skin, rapid heart rate, and extreme weakness.
\end{itemize}

\section*{Diagnosis and Investigations}
The diagnosis is established prior to surgery through careful clinical assessment, and investigations are focused on ensuring patient safety and confirming non-pregnancy:
\begin{itemize}
    \item \textbf{Pre-operative assessment:} Detailed medical history, examination, review of previous abdominal surgeries, and assessment of anaesthetic risk.
    \item \textbf{Pregnancy test:} A urinary or serum beta-hCG test performed on the morning of surgery to rule out an early luteal-phase pregnancy.
    \item \textbf{Cervical screening:} Verification of an up-to-date Cervical Screening Test (CST) as part of routine gynaecological care.
    \item \textbf{Routine blood tests:} Full blood count (FBC) and group and screen may be performed for patients with a history of anaemia or bleeding disorders.
    \item \textbf{Pelvic ultrasound:} Recommended only if there is a suspected pelvic mass, severe endometriosis, or anatomical anomalies detected on physical examination.
\end{itemize}

\section*{Management}
Management involves pre-operative counseling, surgical execution, and post-operative recovery support:
\begin{itemize}
    \item \textbf{Pre-operative counseling:} Detailed discussion regarding the permanence of the procedure, risk of regret, alternatives (such as vasectomy and LARCs), and failure rates.
    \item \textbf{Anaesthesia:} Standard general anaesthesia with endotracheal intubation to protect the airway during abdominal insufflation.
    \item \textbf{Surgical procedure:} Insertion of ports, inspection of the pelvis, and occlusion of both tubes using Filshie clips or performing bilateral salpingectomy.
    \item \textbf{Analgesia:} Multimodal pain relief including intra-operative local anaesthetic infiltration at incision sites, post-operative paracetamol, and non-steroidal anti-inflammatory drugs (NSAIDs).
    \item \textbf{Activity modification:} Rest and avoidance of heavy lifting (greater than 5 kg) or strenuous exercise for approximately 1 to 2 weeks.
    \item \textbf{Follow-up:} A post-operative check-up with the gynaecologist or general practitioner at 2 to 6 weeks to review wound healing and discuss histopathology if a salpingectomy was performed.
\end{itemize}

\section*{Complications}
The potential complications of laparoscopic sterilisation include:
\begin{itemize}
    \item \textbf{Surgical injury:} Accidental damage to the bowel, bladder, ureters, or major blood vessels (e.g. aorta or iliac vessels) during trocar insertion.
    \item \textbf{Anaesthetic risks:} Adverse drug reactions, aspiration pneumonia, or cardiovascular instability.
    \item \textbf{Infection:} Superficial wound infection, pelvic abscess, or peritonitis.
    \item \textbf{Failure of sterilisation:} A small risk of pregnancy (approximately 1 in 200 to 500 women over their lifetime) due to tubal recanalisation or clip slippage.
    \item \textbf{Ectopic pregnancy:} If pregnancy does occur, there is a significantly high chance that it will be ectopic, which is a medical emergency.
    \item \textbf{Psychological regret:} Long-term regret, particularly if the procedure was performed at a young age or under pressure from a partner.
\end{itemize}

\section*{Prognosis and Outcomes}
The prognosis for patients undergoing laparoscopic sterilisation is excellent. Most women return home on the same day as the surgery and resume light, non-strenuous daily activities within 2 to 7 days. Complete wound healing and full recovery occur within 2 weeks. The contraceptive efficacy is immediate (unlike vasectomy) and extremely high, with a lifetime failure rate of less than 0.5\%. The procedure is permanent, and while surgical reversal is technically possible in some cases of tubal occlusion, it is complex, expensive, and has low success rates. Bilateral salpingectomy is completely irreversible but provides the added benefit of significantly reducing the risk of ovarian cancer.

\section*{Prevention and Self-Care}
Self-care measures to support recovery and prevent complications include:
\begin{itemize}
    \item \textbf{Wound hygiene:} Keep the incision sites clean and dry. Avoid soaking in baths, spas, or swimming pools until the wounds are fully healed.
    \item \textbf{Shoulder pain management:} Lie flat or use heat packs on the shoulder to alleviate discomfort from trapped carbon dioxide gas.
    \item \textbf{Stool softeners:} Use mild laxatives or eat high-fibre foods to prevent straining during bowel movements, which can aggravate abdominal pain.
    \item \textbf{Pelvic rest:} Avoid vaginal intercourse and the use of tampons for at least 1 to 2 weeks to reduce the risk of pelvic infection.
    \item \textbf{Graduated activity:} Walk gently to promote circulation and prevent deep vein thrombosis, but avoid vigorous exercise.
\end{itemize}

\section*{Sources and Further Reading}
The following reputable organisations provide further information on laparoscopic sterilisation:
\begin{itemize}
    \item Royal Australian and New Zealand College of Obstetricians and Gynaecologists (RANZCOG). \textit{Patient information on female sterilisation.} https://ranzcog.edu.au/
    \item Healthdirect Australia. \textit{Overview of female sterilisation procedures.} https://www.healthdirect.gov.au/
    \item NHS. \textit{Guide to female sterilisation.} https://www.nhs.uk/
    \item Family Planning Victoria / Sexual Health Victoria. \textit{Information on permanent contraception.} https://www.shvic.org.au/
    \item Mayo Clinic. \textit{Details of tubal ligation and laparoscopic approaches.} https://www.mayoclinic.org/
\end{itemize}
